
\begin{table}[h!]
\centering
\renewcommand{\arraystretch}{1.1}
\resizebox{0.95\linewidth}{!}{\begin{tabular}{llcc}
\toprule
\textbf{Module} & \textbf{Operation} & \textbf{Auto-Regressive LM (decode)} & \textbf{Diffusion LM (length $L$)} \\
\midrule
\multirow{4}{*}{MHA} 
& $W_Q$, $W_K$, $W_V$ projections & $\mathcal{O}(d^2)$ & $\mathcal{O}(Ld^2)$ \\
& Query $\times$ Key ($QK^\top$) & $\mathcal{O}(l^2 d / h)$ & $\mathcal{O}(L^2 d / h)$ \\
% & Softmax + scaling              & $\mathcal{O}(l^2)$        & $\mathcal{O}(L^2)$ \\
& Attention Score $\times$ Value           & $\mathcal{O}(l d / h)$ & $\mathcal{O}(L^2 d / h)$ \\
& Output projection ($W_{\text{out}}$) & $\mathcal{O}(d^2)$ & $\mathcal{O}(L d^2)$ \\
\midrule
\multirow{2}{*}{FFN}
& $W_1$ projection               & $\mathcal{O}(d \dot d_{\text{ff}})$ & $\mathcal{O}(L d d_{\text{ff}})$ \\
& $W_2$ projection               & $\mathcal{O}(d d_{\text{ff}})$ & $\mathcal{O}(L d d_{\text{ff}})$ \\
\bottomrule
\end{tabular}}
\vspace{1.0em}
\caption{Compute complexity of Transformer modules: AR decodes a prefix of length \(l\) per token vs. DLM processes full length \(L\) each denoising step.}
\label{tab:transformer-complexity}
\end{table}

\subsubsection{Bottleneck of Diffusion Language Model Inference}

While diffusion models allow parallelized generation of tokens, this comes at the expense of extra computation overhead per forward pass. 
Given the input $x_T$ with sequence length $L$, each forward pass of the diffusion model $f_\theta$ requires full-sized computation~(both current and past tokens) throughout the multi-head attention~(MHA) block and feed-forward network~(FFN) of the transformer model. 
Unlike the autoregressive model where the past tokens are cached for the new token generation, diffusion model produces output without the consideration of causality. 
As shown in Table~\ref{tab:transformer-complexity}, DLM models introduce an additional computational complexity of $O(L)$ at each module of a transformer-based model architecture. 
As a result, the latency overhead introduced by the DLM architecture is massive compared to the Auto-Regressive Model, as shown in the comparison in Section~\ref{sec:res}. 
Based on the theoretical analysis above, the first challenge of the diffusion model is: 

\textbf{Challenge}\circled{1}: \textit{How to alleviate the latency and computation overhead caused by the liability of full sized cache computation of DLM?}


% it incur an additional factor of \(O(L)\) cost per step from recomputing full‐sequence forward passes. 

% In addition, quality drops significantly as the size of the unmasked set $U_t$ increases. In this work, we address this inefficiency with two complementary training-free techniques, KV approximation caching \zhanqiu{Rename} and AR-guided unmasking, which accelerate inference without compromising quality.  

% Let L denote the full sequence length, which is equal to the sum of input token length and context window length for generation tasks. Each $f_\theta$ is a full forward pass of a transformer decoder with input sequence length of L, which includes multi-head attention (MHA) and feedforward (FFN) layers, applied to the entire sequence at every timestep.

% Table~\ref{tab:transformer-complexity} highlights a key contrast in forward pass efficiency between autoregressive models and diffusion models. Autoregressive decoding operates on a growing prefix of length $l$, while diffusion models operate on the full sequence of length $L$ at each denoising step. All linear projections in DLM scale with $L$ compared to the constant size in ARM. 

% Most linear projections in ARM involve matrix-vector multiplication (e.g., $\mathbb{R}^{d \times d} \cdot \mathbb{R}^{d \times 1}$), while DLM requires matrix-matrix multiplication ($\mathbb{R}^{d \times d} \cdot \mathbb{R}^{d \times L}$), leading to significantly higher computational cost.
